% Generated draft: requires mathematical and source review.
\documentclass[11pt]{article}
\usepackage[margin=25mm]{geometry}
\usepackage{fontspec}
\setmainfont{Latin Modern Roman}
\usepackage{amsmath,amssymb,mathtools,unicode-math}
\setmathfont{Latin Modern Math}
\usepackage{enumitem,needspace,xurl,hyperref}
\setlength{\parindent}{0pt}
\setlength{\parskip}{5pt}
\setlength{\emergencystretch}{3em}
\newcommand{\N}{\mathbb{N}}
\newcommand{\Z}{\mathbb{Z}}
\newcommand{\Q}{\mathbb{Q}}
\newcommand{\R}{\mathbb{R}}
\newcommand{\C}{\mathbb{C}}
\newcommand{\F}{\mathbb{F}}
\newcommand{\A}{\mathbb{A}}
\newcommand{\PP}{\mathbb{P}}
\newcommand{\E}{\mathbb{E}}
\newcommand{\eps}{\varepsilon}
\newcommand{\dd}{\,\mathrm{d}}
\begin{document}
\section*{UnsolvedMath ID: 30006305}
\textbf{Source problem code:} catalogue entry
\subsection{Definitions and mathematical statement}
A toric variety is a normal variety containing an algebraic torus as a dense open subset, with the torus multiplication extending to an action on the variety. Two integral varieties are birational if they have isomorphic dense open subsets. For a variety \(V\) and a closed center \(C\subset V\) with ideal sheaf \(\mathcal I_C\), the ordinary blowup is
\[
\operatorname{Bl}_C(V)=\operatorname{Proj}_V\!\left(\bigoplus_{m\geq 0}\mathcal I_C^m\right).
\]
Here a smooth blowup is interpreted as an ordinary blowup along a smooth proper closed center; a blowdown is the reverse step along its blowup morphism.

The supplied question asks whether every pair of smooth birational toric varieties \(X,Y\) admits a diagram
\[
X=X_0\dashrightarrow X_1\dashrightarrow\cdots\dashrightarrow X_r=Z=Y_s\longrightarrow Y_{s-1}\longrightarrow\cdots\longrightarrow Y_0=Y,
\]
where \(r,s\in\Z_{\geq 0}\), all displayed varieties are smooth, and
\[
X_i\cong\operatorname{Bl}_{C_i}(X_{i-1})\quad(1\leq i\leq r),\qquad
Y_j\cong\operatorname{Bl}_{D_j}(Y_{j-1})\quad(1\leq j\leq s)
\]
for smooth proper closed centers \(C_i\subset X_{i-1}\) and \(D_j\subset Y_{j-1}\). Thus all blowups occur before all blowdowns. The arrows from \(X_{i-1}\) to \(X_i\) denote the rational inverses of the blowup morphisms.

\emph{Unresolved scope.} This diagram makes the order of operations explicit, but the supplied record does not determine the full hypotheses. It does not specify the base field, characteristic, splitness of the tori, completeness or projectivity, whether the intermediate varieties must be toric and the centers torus invariant, or whether a prescribed torus-equivariant birational map must be factored. If a map \(\phi\colon X\dashrightarrow Y\) is prescribed, the additional requirement is
\[
\phi=q\circ p^{-1},
\]
where \(p\colon Z\to X\) and \(q\colon Z\to Y\) are the composite blowup morphisms. These additional requirements are not asserted by the supplied clean statement. The fragmentary original text does not resolve them, so the diagram is a provisional rendering of the supplied question rather than a fully identified formulation of Oda's conjecture.
\subsection{Short English statement}
Can every pair of smooth birational toric varieties be connected by first performing finitely many ordinary blowups along smooth centers and then finitely many blowdowns of that kind? The supplied record leaves the base field, completeness assumptions, toric equivariance, and any requirement to factor a specified map unresolved.
\subsection{Sources}
\begin{itemize}
\item Local UnsolvedMath registry, permanent ID \(30006305\), exact source problem code \texttt{catalogue entry}, title: ``Oda's Strong Factorization Conjecture for Toric Varieties.'' The record belongs to the Oberwolfach Reports Open Problems collection and assigns the year 2025.
\item \emph{Toric Geometry} (2025), Oberwolfach Reports, pp.\ 883--940, DOI \texttt{10.4171/owr/2025/19}: \url{https://doi.org/10.4171/owr/2025/19}. The supplied original fragment reads: ``What can be said about their relation? Oda's strong factorization conjecture (and relatives) Karim Adiprasito.'' Karim Adiprasito is named in that fragment; the registry supplies no separate proposer attribution.
\item Adiprasito and Pak, \emph{All triangulations have a common stellar subdivision}: \url{https://arxiv.org/abs/2404.05930}.
\item Da Silva and Karu, \emph{On Oda's Strong Factorization Conjecture}: \url{https://doi.org/10.2748/tmj/1309952085}.
\item Supplied literature assessment dated exactly \texttt{2026-08-22}: the registry labels the problem open and describes the classical smooth-center statement as remaining open, while reporting that a weighted or rationally smooth version was proved. This is supplied dated curation, not an independently verified literature assessment; the precise alternative theorem and its hypotheses are not supplied.
\item The registry labels the statement \texttt{corrected\_verified} and reports that the official report was checked and the question isolated and repaired for self-contained publication. These are supplied verification claims, not verification performed for this draft.
\end{itemize}
\end{document}
